% Source: https://artofproblemsolving.com/community/c3h1993731_znzx
% Stable identity: aops-topic-1993731-q01
(a) If $ab\equiv1\pmod n$, then $ab-1=kn$ for some integer $k$, so any common divisor of $a$ and $n$ divides $ab-kn=1$; hence $\gcd(a,n)=1$. Conversely, if $\gcd(a,n)=1$, Bezout's identity gives integers $x,y$ with $ax+ny=1$, so $ax\equiv1\pmod n$ and $x$ is an inverse.

(b) By (a) the elements are the classes of the integers in $\{1,\dots,14\}$ coprime to $15$:
\[
1,\ 2,\ 4,\ 7,\ 8,\ 11,\ 13,\ 14 .
\]
Their inverses: $1\cdot1=1$; $2\cdot8=16\equiv1$; $4\cdot4=16\equiv1$; $7\cdot13=91\equiv1$; $11\cdot11=121\equiv1$; $14\cdot14=196\equiv1$. So $1,4,11,14$ are their own inverses, while $2\leftrightarrow8$ and $7\leftrightarrow13$ are inverse pairs. There are $\boxed{8}$ elements, which equals $\varphi(15)=15\left(1-\dfrac{1}{3}\right)\left(1-\dfrac{1}{5}\right)=8$. (The ``$\times$'' in the notation means ``multiplicative''; with multiplication modulo $n$ these classes form a group.)
